% Example content adapted from David Molik’s tenure document.
\PortfolioSection{List of Candidate’s Research, Scholarly and Creative Activities, and Discovery (RSCAD)}
\SectionGuidance{Present your research and creative work, including publications, presentations, and other relevant outputs. State the status of each work and your contribution to collaborative activities.}
\ExampleNote

During the evaluation period, I engaged in sustained research and creative activity directly related to my directed service responsibilities in scholarly communication, publishing infrastructure, research data governance, and cyberbiosecurity. Although I am in the early stage of my appointment at Kansas State University, I have completed peer-reviewed publications, delivered scholarly presentations, and contributed to grant-funded research coordination activities that demonstrate an active and coherent research agenda.

\subsection*{Peer-Reviewed Publications}

I co-authored three peer-reviewed journal articles published in 2025 that contribute tools, workflows, and standards supporting research infrastructure and FAIR data practices:

\begin{itemize}
  \item \textit{AgBioDatabase Finder: an online tool to help researchers find and submit agricultural genomic, genetic, and breeding data} (\textit{microPublication Biology}, 2025), which presents a community-facing tool to improve discoverability and submission of agricultural genomics data.

  \item \textit{BALROG-MON: a high-throughput pipeline for Bacterial AntimicrobiaL Resistance annOtation of Genomes--Metagenomic Oxford Nanopore} (\textit{microPublication Biology}, 2025), documenting a reproducible computational pipeline supporting antimicrobial resistance annotation.

  \item \textit{Guidelines for Gene and Genome Assembly Nomenclature} (\textit{Genetics}, 2025), a community standards article establishing best practices for gene and genome assembly nomenclature to support interoperability and data reuse.
\end{itemize}

In 2026, this trajectory continued with \textit{Toward standardization in arthropod and biodiversity genome projects} (\textit{Genetics}, 2026), a community standards article reviewing reporting standards across arthropod genome projects to support interoperability and reuse of biodiversity genomic data.

While portions of these projects were initiated prior to my appointment, substantive contributions, final revisions, and publication support spanned the appointment transition. In particular, BALROG-MON is profile-dated February 19, 2025, days before the February 25, 2025 appointment recorded in the portfolio. Collectively, these publications demonstrate applied research and creative activity aligned with my role in scholarly communication and research infrastructure.

\subsection*{Scholarly Presentations and Community Dissemination}

I contributed to scholarly dissemination and community-building through conference organization and invited presentations. I organized the session \textit{Cyberbiosecurity in Agricultural Genomics} and co-organized the session \textit{The AgBioData Consortium: Challenges and Recommendations for FAIR Genetic, Genomic, and Breeding Data} at the Plant and Animal Genome Conference (PAG33).

In January 2026, I delivered the scholarly talk \textit{Some Thoughts on Genomics Data Integrity} at the \textit{Cyberbiosecurity in Agricultural Genomics} workshop at PAG33. This presentation articulated emerging challenges related to data integrity, provenance, and trust in genomics and reflected ongoing conceptual work that informs both my research agenda and my directed service responsibilities.

\subsection*{Grant-Funded and Collaborative Research Activity}

I served as a Co-Principal Investigator on the National Science Foundation Research Coordination Network project \textit{``Reimagining a Sustainable Data Network to Accelerate Agricultural Research and Discovery''} (Award \#2126334). Through this role, I contributed to grant-funded coordination activities, including annual meetings, working groups, and community assessment efforts focused on data standards, interoperability, and scholarly communication in agricultural genomics.

This grant-supported work directly informs my research and creative activity by providing empirical grounding and community context for future publications and presentations related to research data governance and infrastructure.

\subsection*{Research Trajectory}

Together, these research and creative activities articulate a coherent scholarly trajectory centered on the design, governance, and sustainability of open scholarly communication and research infrastructure. Rather than pursuing domain science in isolation, my research focuses on the systems, standards, and practices that enable research to be created, disseminated, preserved, and reused at scale. This emphasis is directly aligned with the mission of the Open Publishing Exchange (OPE), which integrates scholarly publishing, open education, and research data services within the Libraries.

My peer-reviewed publications, standards work, and conference presentations contribute empirical tools, reproducible workflows, and governance frameworks that directly inform OPE services, including open publishing platforms, data curation initiatives, and institutional research support. For example, work on FAIR data practices, genomic data nomenclature, and data integrity provides a research-informed foundation for how OPE approaches repository services, publishing workflows, and policy guidance for faculty and researchers.

By grounding departmental service development in active research and community standards leadership, this scholarly agenda strengthens the Libraries’ ability to respond to evolving funder mandates (e.g., The Nelson Memo), disciplinary norms, and technological change. In this way, my research trajectory is not ancillary to my administrative role but is integral to advancing OPE as a research-facing, standards-informed service organization. This integration provides a sustainable foundation for continued research productivity that simultaneously advances institutional infrastructure and contributes to national and international scholarly communities.

\Redacted{Additional working updates}{Private consultation, proposal, and working-record details are omitted from this example.}

\section{Research and Creative Evidence}
\SectionGuidance{Provide publication records, persistent links, and other appropriate evidence for the research and creative work described above.}
\printbibliography[keyword={creative},heading=none]
